\subsection{Kuhn Poker}
Die simpelste definierte Pokervariante ist One Card Poker auch Kuhn Poker gennant.
Der Name Kuhn Poker stammt durch den Erfinder dieses Spiels Harald Kuhn, welcher es in seinem Paper aus dem Jahr 1951 mit dem Titel \grqq Simplified Two-Person Poker \glqq beschrieb. \\
Kuhn Poker wird mit einem Kartendeck bestehend aus 3 Karten gespielt.
In seinem Paper sind die Karten mit 1,2 und 3 bennant, wobei 1 die schlechteste und 3 die beste Karte ist.
Für diese Arbeit wird die Notation J, Q, K verwendet, die den Rängen eines Standard-Kartenspiels entspricht.

Zu Beginn zahlen beide Spieler einen Pflichteinsatz in Höhe eines Chips in den Pot, dies wird auch Ante genannt. 
Beide Spieler erhalten eine private Karte und es werden keine öffentlichen Karten ausgeteilt.
Nach dem austeilen der privaten Karte startet die erste und einzige Setzrunde, in welcher die Spieler jeweils einen Chip in den Pot setzen dürfen, womit das Setzlimit bei Eins liegt.
Im Falle eines Showdowns, gewinnt die höhere Karte $J < Q < K$.

Kuhn definiert in seinem Paper auch noch andere Versionen seines Spiels indem er das Verhältnis zwischen Ante und Setzgröße anpasst.
Der eben beschriebene Fall in welchem der Ante genauso groß ist wie die Setzgröße, wird von ihm Fall 2 genannt.
Fall 1 liegt vor, wenn die Setzgröße kleiner als der Ante ist.
Fall 3 tritt ein, wenn der Ante kleiner als die Setzgröße ist und die Setzgröße kleiner als das Doppelte des Antes ist.
Fall 4 liegt vor, wenn die Setzgröße dem Doppelten des Antes entspricht.
\cite{kuhn1951simplified}